\begin{exercise}{76}
Show that the intersection of left- and right-anodyne maps strictly contains the
inner-anodyne maps.

\begin{proof}
Every left or right fibration is an inner fibration, so every
inner-anodyne map is both left- and right-anodyne
(Definitions~1.3.1 and~1.3.2).
We show that the vertex inclusion $\{0\}\to J$ belongs to both
classes; Exercise~\ref{ex:75} shows that it is not inner-anodyne.

Use the tuple description of $J$ from Exercise~\ref{ex:60}:
its nondegenerate simplices are precisely the alternating
tuples in $0,1$.
Let $\alpha_n=(0,1,0,1,\ldots)$ be the alternating $n$-simplex
starting at $0$, and let $K_n\subset J$ be the simplicial subset
it generates.
For $n\geq1$, $K_n$ contains both alternating simplices in
each dimension below $n$, and only $\alpha_n$ in dimension $n$:
the lower-dimensional ones occur as consecutive segments
starting at position $0$ or $1$.
Consequently $K_0=\{0\}$, $K_{n-1}\subset K_n$, and
$\bigcup_nK_n=J$.

The map $\tau_n\colon\Delta^n\to J$ classifying $\alpha_n$
gives a pushout square for every $n\geq1$:
\[
\begin{tikzcd}
 \Lambda_0^n \arrow[r,hook] \arrow[d] & \Delta^n \arrow[d,"\tau_n"] \\
 K_{n-1} \arrow[r,hook] & K_n.
\end{tikzcd}
\]
Indeed, $d_n\alpha_n=\alpha_{n-1}$, while deleting an interior
entry creates two adjacent equal entries and hence a degeneracy
of a lower-dimensional alternating simplex.
Thus all faces other than $d_0\alpha_n$ lie in $K_{n-1}$.
To check the pushout degreewise, a $k$-simplex
$r\colon[k]\to[n]$ outside $\Lambda_0^n$ visits every vertex
$1,\ldots,n$.
Its image is therefore either $[n]$ or $\{1,\ldots,n\}$.
In the first case $\tau_n(r)$ is a degeneracy of $\alpha_n$;
in the second it is a degeneracy of the alternating
$(n-1)$-simplex starting at $1$.
These are exactly the simplices of $K_n$ absent from $K_{n-1}$.
In each case their constant runs recover $r$ uniquely, and
every such degeneracy occurs.
Hence the complement of the horn maps bijectively onto
$K_n\setminus K_{n-1}$ in every degree, proving the claim,
including $n=1$.
By Corollary~1.3.16, the resulting countable composition
$\{0\}\to J$ is left-anodyne.

Taking opposites preserves these pushouts and their colimit,
and carries $\Lambda_0^n\to\Delta^n$ to
$\Lambda_n^n\to\Delta^n$ by Exercise~\ref{ex:21}.
Thus $\{0\}\to J^{\mathrm{op}}$ is right-anodyne.
Inversion in the groupoid defining $J$ induces
$J^{\mathrm{op}}\cong J$, fixing both objects
(Exercise~\ref{ex:21}(1)).
In particular it fixes the chosen vertex $0$, so the same
inclusion $\{0\}\to J$ is right-anodyne as well.
\end{proof}
\end{exercise}
